% 本文件由 LaTeX 排版 Agent 根据有效 translation.json 自动生成。
% author=Sulpitius Severus (c. 363-c. 420); work_id=3504; title=可疑信件

\chapter{\WorkTitle{可疑信件}}

% structure: p0001 source=h1 target=omitted reason=page-title-already-rendered-by-parent

% structure: p0002 source=h2 target=section reason=relative-heading-rank; base=h2

\section{\WorkTitle{\Person{苏尔皮奇乌斯·塞维鲁}的可疑信件}}

% structure: p0003 source=h2 target=section reason=relative-heading-rank; base=h2

\section{\Emph{第一封信}}

\Emph{圣司铎\Person{塞维鲁}关于最后审判致其妹\Person{克劳迪娅}的一封信}\par

% structure: p0005 source=h3 target=subsection reason=relative-heading-rank; base=h2

\subsection{\Emph{第一章}}

拜读你的来信，我的心情在诸多方面深受触动，无法抑制泪水。因为我既因喜乐而哭泣——从你信件的话语中，我能看出你正按着主天主的诫命生活——又因对你的极度思念而不得不哀叹，我无过却与你分离；若非你寄信给我，这种感觉会更加强烈。难道我不该与这样的妹妹相伴吗？但我以你的救恩为证，我常渴望到你那里去，但至今仍被那惯于阻碍我们者的反对所阻。因我热切渴望，既急于通过见你满足心愿，又似乎觉得，我们若相遇，能更有效地成就主的工作，因彼此安慰而活在这重担压下的世界之上。但我如今不定下探访你的日子或时间，因每当我如此行，总不能实现目的。我将等候主的旨意，并希望藉我的祈求与你的祈祷，祂能成就此事，使我们从恒心中获益。\par

% structure: p0007 source=h3 target=subsection reason=relative-heading-rank; base=h2

\subsection{\Emph{第十二章}}

然而，因你曾渴望从我寄给你的所有信件中获得滋养你生命与信德的训导，以致因我频繁写信给你，如今已穷尽此类言辞；实在无法再写任何新东西给你，以免重复先前所写。而因天主的恩宠，你如今已无需被劝勉，因你在圣洁生命的开端便圆满信德，彰显在基督内的热诚之爱。但我仍敦促你一事：不可退回已离弃的事物，不可再贪恋已蔑视之物，既已扶犁，不可再向后看\BibleRef{路 9:62}，倒退脚步；因为若陷入此过，你的犁沟必失其直，耕耘者将不得其应得的酬报。况且，若稍有缺失，连一份酬报也无法获得。因为正如我们必须从罪恶逃向义德，同样，已进入义德实践的人必须提防，以免使自己暴露于罪恶。经上记载：\BibleQuote{当义人转离其义而行恶时，他的义不再被记念。}\BibleRef{则 18:24}为此，我们必须站立，必须劳苦，使已逃离罪恶者不致失去预备好的酬报。因为仇敌伺机攻击我们，要一击打倒那已卸下信德之盾的人。因此，我们的盾牌不可丢弃，免得肋旁暴露于攻击；我们的剑也不可收起，免得仇敌开始无所畏惧。而且我们知道，若他见一人全副武装，便会退却。我们并非不知，天天与肉躯及世界争战是艰难的事。但若你思念永恒，思念天国——主必屈尊赐予我们这些罪人——请问，有何苦难足够巨大，以致我们能配得这些？此外，我们在世上的奋斗是短暂的；因为即使死亡不来速临，衰老也会到来。岁月流逝，光阴飞驰；而我盼望主耶稣必速速呼召我们到祂身边，因我们为祂所爱。\par

% structure: p0009 source=h3 target=subsection reason=relative-heading-rank; base=h2

\subsection{\Emph{第三章}}

哦，我们的离世何等有福，那时基督将接我们进入祂的居所，我们已藉更美好生命的经验，从罪恶的污秽中被炼净！殉道者与先知将迎接我们，宗徒将与我们同行，天使将欢欣，总领天使将喜乐，而撒殚被征服后脸色苍白，却仍保留其残酷面目，因他失去从我们罪恶中为自己获取的一切利益。他将看见藉慈悲赐予我们的光荣，功绩因光荣而受尊崇。我们将战胜被征服的敌人。那时，世上的智者何在？贪得无厌者、奸淫者、不敬者、酗酒者、恶言者将在何处被认出？这些可悲之人将如何为自己辩护？\Quote{吾主，我们不曾认识祢；我们未见祢在世；祢未曾派遣先知；祢未赐法律给世界；我们未见列祖；我们未读圣贤生平。祢的基督从未在地上；伯多禄沉默；保禄拒绝宣讲；无福音者教导。未有殉道者可作典范；无人预言祢将来的审判；无人命我们给穷人穿衣；无人禁止我们节制情欲；无人劝我们对抗贪念；我们因无知而跌倒，不知自己所作何事。}\par

% structure: p0011 source=h3 target=subsection reason=relative-heading-rank; base=h2

\subsection{\Emph{第四章}}

反对他们，圣人群中的义人诺厄将首先宣告：\Quote{我，主，曾预言因人的罪恶将有洪水来临，洪水之后，我以自身为善人树立榜样；因为我未与丧亡的恶人一同丧亡，使他们知道无辜者的救恩与罪人的惩罚是什么。}此后，忠信的亚巴郎将反驳他们说：\Quote{我，主，约在世界年代的中期，奠定了信仰的基础，使人性因你而相信；我被选为万民之父，使他们效法我的榜样；我毫不迟疑，主，将年少的依撒格献给你为祭献，使他们明白，当他们见到我连独生子也不顾惜时，就应知道没有什么是不应献给主的；我，主，遵照你的命令离开了我的故乡和家族，为使他们也有一个榜样，教导他们离开世界和时代的邪恶；我，主，是第一个在肉身形态下认出你的人，虽你以不同于你自己的形态向我显现，我仍毫不犹豫地相信我所看见的是谁，为使他们学会不按肉身、而按精神判断。}有福的梅瑟将在他恳求中支持他，说：\Quote{我，主，按你的命令将法律交给了所有这些，为使那些自由信仰不能影响的人，至少口头上的法律能约束他们：我说：\InnerQuote{\BibleQuote{不可奸淫} \BibleRef{出 20:14}}，为的是阻止淫乱放荡；我说：\InnerQuote{\BibleQuote{应爱你的近人} \BibleRef{肋 19:18}}，为的是增加爱德；我说：\InnerQuote{\BibleQuote{你只应崇拜主} \BibleRef{申 6:13}}，为的是他们不向偶像献祭，或容许庙宇存在；我命令不可作假见证，为的是封住这些人的口，反对一切谎言。我借着你的大能之灵在我内的运行，陈述了从世界起初所成就和所说的事，为使过去的知识能将这些人的教训传达到未来的事。我预言了，主耶稣，你的来临，为使这来临对这些人不至出乎意料，当他们被召去承认我所先前宣布要来的那一位时。}\par

% structure: p0013 source=h3 target=subsection reason=relative-heading-rank; base=h2

\subsection{第五章}

此后，与主后裔相称的达味将站起来宣告：\Quote{我，主，用各种方式宣扬了你；我阐明惟独你的名应受敬拜；我说：\InnerQuote{敬畏主的人是有福的}；我又说：\InnerQuote{圣者必在光荣中欢乐}；我还说：\InnerQuote{恶人的欲望必归乌有}，为使这些人认识你，停止犯罪。我，当获得王权时，身穿麻衣，铺上灰尘，放下我伟大的标志，躺卧在衣服中，为给这些人谦逊和柔和的榜样。我宽恕了那些想要杀我的仇敌，为使他们赞同我的慈悲，视为值得效法的。}此后，堪当天主圣神的依撒意亚不会沉默，而要说：\Quote{我，主，当你藉我的口说话时，我发出警告：\InnerQuote{\BibleQuote{祸哉，那些房屋连房屋的人} \BibleRef{依 5:8}}，为的是给贪得无厌设限。我作证你的忿怒临到恶人，为的使这些人至少因惧怕惩罚，而非希望赏报，远离恶行。}\par

% structure: p0015 source=h3 target=subsection reason=relative-heading-rank; base=h2

\subsection{第六章}

此后，以及其他几位为我们完成教导职责的之后，\Person{天主子}亲自将这样说：\Quote{我，当然，高居尊位，手握上天，拳握大地，内外伸展，在一切受造物的内部，在一切运动事物的外部，不可捉摸，本性之能无限，目不能见，触不能及，为了成为你们中最微小者（为制服你们心灵的刚硬，藉健全教义软化你们的不信），我屈尊降生成人，并放下天主的荣耀，取了奴仆的形体，为能同你们分享肉身的软弱，进而因服从救恩的诫命，带你们分享我的荣耀。我恢复病人和虚弱者的健康，聋子的听觉，瞎子的视觉，哑巴的言语能力，瘸子的行走能力；为藉天上的征兆，更容易影响你们相信我和我所宣告的事，我应许你们天国；我还为使你们有逃避刑罚的榜样，将几乎在死亡时刻认我的强盗置于乐园，为使你们甚至追随那被认为获得赦罪者的信德。并且，我为你们树立榜样，使你们也能受苦；我为你们受苦，为使无人犹豫为自己受天主为人所受的苦。我复活后显现，为不使你们的信仰倾覆。我在\Person{伯多禄}身上劝诫犹太人；我在\Person{保禄}身上向外邦人宣讲；我不后悔这样做，因为美果随之而来。善人明白我的工作；忠信者成全了它；义人完成了它；慈悯者成就了它：有大量殉道者，大量圣人。我提到的这些人无疑与你们在同一肉身和同一世界。那么，为何我在你们身上找不到善工，你们这些毒蛇的孽子？你们甚至在地世生命的末刻，也不为自己的恶行悔改。你们用嘴唇尊敬我，却用行为否认我，这有什么益处呢？如今你们的财富、尊荣、权势和享乐在哪里？我向你们宣布的不再是新判决：你们只是承受我从前预言过的审判。}\par

% structure: p0017 source=h3 target=subsection reason=relative-heading-rank; base=h2

\subsection{第七章}

然后，圣史将向那些可怜的人重复这句话：\BibleQuote{去外方的黑暗里，那里要有哀号和切齿。\BibleRef{玛 22:13}} 你们这些可怜的人啊，这些话如今竟不触动你们！那时他们将看见自己的惩罚，看见他人的荣耀。让他们享用今世吧，只要他们不享有那为圣人预备的永恒。让他们富足于财宝，安卧于黄金；只要在那里他们被发现是匮乏穷困的。让他们在今世富有，只要他们在永恒中贫穷，因为经上论及他们记载：\Quote{富足者有所缺乏，忍受饥饿。} 但圣经又论及善人添加了以下内容：\Quote{寻求主的人，必不缺乏一切美物。}\par

因此，我的姊妹，尽管那些人讥讽我们，尽管他们称我们为愚昧可悲之人，我们却更应因这些辱骂而喜乐欢欣，因为这些辱骂为我们积累荣耀，为他们积累惩罚。我们不要嘲笑他们的愚昧，反而要为他们不幸而悲伤；因为他们中有许多我们自己的人，若我们能争取他们，我们的荣耀必将增加。但无论他们怎样行事，让他们成为我们眼中的外邦人和税吏；我们自己却要保持平安妥当。若他们如今因我们哀哭而欢喜，日后必将轮到我们因他们的痛苦而欢喜。祝你平安，最亲爱的姊妹，在基督内深蒙爱怜。\par

% structure: p0020 source=h2 target=section reason=relative-heading-rank; base=h2

\section{第二封信}

苏尔皮基乌斯·塞维鲁致其姊妹克劳迪娅论童贞书\par

% structure: p0022 source=h3 target=subsection reason=relative-heading-rank; base=h2

\subsection{第一章}

圣洁的童贞在天上的恩赐中，享有何等大的福乐，除了圣经的见证外，我们也从教会的实践中得知，借此我们被教导，那些藉特别奉献而献身于童贞者，拥有独特的功绩。因为虽然全体信众接受同等的恩赐，并喜乐于圣事相同的祝福，但童贞者却拥有超越其余人的某些事物，因为他们是从教会圣洁无瑕的团体中被圣神拣选，并由主教在天主祭台上呈献，犹如更圣洁更纯洁的祭献，因他们自愿奉献的功绩。这实在是一个堪当配得上天主的祭献，因为是如此宝贵之存有的献祭，而没有什么比献上祂自己的肖像更能取悦祂。因为我认为宗徒特别指这类祭献，当他说：\BibleQuote{所以弟兄们，我以天主的仁慈请求你们，将你们的身体献上，当作生活的、圣洁的、蒙悦纳的祭献给天主。\BibleRef{罗 12:1}} 因此，童贞拥有别人有的和没有的；既获得普通的恩宠，也获得特殊的恩宠，并且（可以这么说）因它独特的奉献特权而欢喜。因为教会的权威允许我们称童贞者为基督的新娘；同时，仿效新娘的方式，它为那些被祝圣给主的童贞者蒙上头纱，公开表明那些逃离血肉亲情的人特别将要拥有属灵的婚姻。那些人按照婚姻的类比，以属灵的方式与天主结合，他们因爱天主而藐视人的结合。在他们身上，宗徒的那句话得到最充分的应验：\BibleQuote{与主结合的，成为一神。\BibleRef{格前 6:17}}\par

% structure: p0024 source=h3 target=subsection reason=relative-heading-rank; base=h2

\subsection{第二章}

因为舍弃奢华，以心灵的刚毅熄灭被青春火炬点燃的贪欲之火，以属灵的努力压制本性的快感，逆生活法则而行，蔑视婚配的安慰，鄙弃来自子女的甜蜜享受，并在未来福乐的希望中，视今世生活的一切利益为无物——这实在是伟大而神圣的事，几乎超越有形的本性。正如我所说，这是一项伟大而可钦的德行，并理应按其劳苦之伟大，获得极大的酬报。经上说：\BibleQuote{我要赐给阉人，主说，在我家中、在我墙内，有比子女更好的地方和名字；我要赐给他们一个永不废去的永久名字。\BibleRef{依 56:5}} 主又在福音中论及这类阉人，说：\BibleQuote{因为有些阉人，是为了天国而自阉的。\BibleRef{玛 19:12}} 与贞洁相系的挣扎实在伟大，但酬报更大；约束是暂时的，酬报却是永恒的。因为蒙福的宗徒若望也论及他们说：\BibleQuote{羔羊无论往哪里去，他们都跟随祂。\BibleRef{默 14:4}} 我认为，这应作如下理解：天上宫廷中没有一处对他们关闭，而神圣居所的一切住所都将向他们敞开。\par

% structure: p0026 source=h3 target=subsection reason=relative-heading-rank; base=h2

\subsection{第三章}

但是，为使童贞的功勋得以更清晰地闪耀，并使人们更深入地理解它如何堪当天主，请思考这一点：主天主、我们的救主，为拯救人类而屈尊就卑取了人性时，他选择了除童贞女的子宫之外别无他处，为要表明这种美德特别令他喜悦；并且，为向两性指明贞洁的幸福，他有一位童贞母亲，而他自己则永远保持相似的状态。他从而在自己的身上为男子，在他母亲的身上为女子，提供了童贞的榜样，借此得以证明，就两性而言，纯洁的幸福状态拥有完整的天主性，因为凡居住在圣子内的，也完全居住在他母亲内。但我为何要费力宣扬贞洁卓越超绝的功勋，并陈明童贞荣耀的美善呢？我并非不知道许多人已就此主题论述过，并以最确凿的理由证明了它的幸福；何况对于任何有思考能力的人而言，一件事越是难以成就，其功勋就越大，这是毋庸置疑的。因为若有人视贞洁为无足轻重或只有微末价值，那必然是他要么对此事无知，要么不愿承担它所带来的辛劳。因此，那些要么不拥有贞洁，要么不情愿地被迫保持贞洁的人，总是贬低贞洁的重要性。\par

% structure: p0028 source=h3 target=subsection reason=relative-heading-rank; base=h2

\subsection{第十四章}

那么，既然我们已经虽然简略地陈述了洁净的困难与功勋，就必须极其谨慎，以免一件本身蕴含大德行、且注定要获得丰厚赏报的事，竟无法结出应有的果实。因为事物越是珍贵，就越被人焦虑地守护。既然有许多事物，除非得到其他事物的帮助，否则无法保持其应有的卓越——例如蜂蜜，除非有蜂蜡的保护和蜂巢的隔间，事实上更准确地说，由这些所支撑，否则就会失去美味，无法独自存在；又如酒，除非存放在气味宜人的容器中，并经常更换沥青，否则就会失去其天然的甘醇之力——所以我们必须万分谨慎，唯恐童贞也可能需要某些必要的辅助，没有这些辅助，它根本无法结出应有的果实，从而这样一件极其困难之事可能毫无益处（尽管人们一直以为它有益），因为它是在没有其他必要附属品的情况下拥有的。因为，除非我理解错误，贞洁之所以得以完整保持，是为了获得天国中的赏报；而十分确定的是，没有人能在不配得永生的情况下获得这赏报。但永生只能通过遵守所有天主的诫命而赢得，圣经为证说：\BibleQuote{你若愿意进入生命，就该遵守诫命。}\BibleRef{玛 19:17}因此，除非一个人遵守了法律的全部规诫，否则无人拥有那生命；而没有这种生命的人，不能成为天国的拥有者，因为在天国里，不是死人，而是活人将统治。因此，童贞虽希望获得天国的光荣，但若不同时拥有那应许永生的东西，仅靠本身毫无益处；借着永生，才拥有天国的赏报。所以，最重要的是，那些完全保持贞洁，并期望从天主的公义获得其赏报的人，必须遵守所有加给我们的诫命，以免保持光荣的贞洁和节欲所付出的辛劳归于徒然。凡熟悉法律的人，无不知道童贞高于诫命或规条，如宗徒所说的：\BibleQuote{论到童贞者，我没有主的命令，但我提出我的意见。}\BibleRef{格前 7:25}因此，当他仅就保持童贞提出意见，而没有颁布命令时，他就承认童贞高于诫命。所以，保持童贞的人，所做的超过诫命的要求。但如果你同时也做了所命令的事，那么你做了超过命令的事才有益处。因为如果你在某一点上做得少，你怎能夸耀自己做得多呢？渴望实现天主的劝谕，就要特别注意遵守诫命；希望获得童贞的赏报，就要紧紧抓住为赢得生命所必需的东西，使你的贞洁能获得酬报。因为正如遵守诫命确保生命，相反，违反则导致死亡。那因不顺从而被判死的人，不能希望得到属于童贞的冠冕；而且，当他真正被交付惩罚时，也不能期望得到对贞洁应许的赏报。\par

% structure: p0030 source=h3 target=subsection reason=relative-heading-rank; base=h2

\subsection{第五章}

如今有三种美德，藉以获得天国的产业。第一是\OriginalTerm{贞洁}{chastity}，第二是\OriginalTerm{轻视世俗}{contempt of the world}，第三是\OriginalTerm{正义}{righteousness}。这三者若结合在一起，对拥有者大有裨益；若分离，则几乎毫无益处，因为每一样都非仅为其本身，而是为着另一件。首先，要求贞洁，以便更容易做到轻视世俗，因为那些不被婚姻束缚的人更容易轻视世俗。其次，要求轻视世俗，以便维持正义，而那些陷于世俗利益和追求俗世享乐的人难以完全保持正义。因此，谁若拥有了第一种美德——贞洁，却没有同时拥有第二种——轻视世俗，那么他几乎白白拥有第一种，因为他没有第二种，而第一种是为第二种而要求的。若有人拥有第一和第二种，却缺乏第三种即正义，他就徒劳无功，因为前两种主要是为第三种而要求的。拥有贞洁以轻视世俗，却没有那为此而拥有前者的东西，又有何益？或者，你为何要轻视世俗之事，若你不遵守正义，而为了正义，你理当拥有贞洁和轻视世俗？因为正如第一种美德是为第二种，第二种是为第三种，同样，第一和第二种是为第三种；若第三种不存在，则第一和第二种都将毫无益处。\par

% structure: p0032 source=h3 target=subsection reason=relative-heading-rank; base=h2

\subsection{第六章}

但或许你在这里说：\Quote{那么，请教给我什么是正义，好让我认识它，从而更容易地完全实践它。}好吧，我将尽我所能简要地为你解释，并使用平常的言辞，因为我们所讨论的主题绝不应因追求华丽的描述而被遮蔽，而应以最简单的表达方式阐明。对于所有人共同必需之事，理应用普通的言语陈述。正义无非就是不犯罪；而不犯罪就是遵守法律的诫命。遵守这些诫命以两种方式维持：即不做任何禁止之事，并努力履行所命令之事。这就是以下语句的含义：\Quote{远离恶，并行善。}因为我不愿你以为正义只在于不作恶，因为不行善也是恶，而且在这两种情况下都违反了法律，因为那说\Quote{远离恶}的，也说了\Quote{并行善。}你若远离恶，却不行善，你就是法律的违反者，而法律不仅因憎恶一切恶行而成就，也因实行善工而成就。确实，你不仅领受了这样的命令：不可剥夺衣服已穿上之人的衣服，反而要用你自己的衣服遮盖那被剥夺的人；也不仅是要取走那有面包之人的面包，而要甘愿将你的面包分给那没有面包的人；也不仅是不可将穷人从他自己的住所赶走，反而要接待那被赶出且无处可归的人进入你的住所。因为我们所领受的诫命是\BibleQuote{与哭泣的人一同哭泣\BibleRef{罗 12:15}}。但我们怎能与他们一同哭泣，如果我们不分享他们的任何需要，不在他们哀叹的事上提供帮助？因为天主并不要求我们无果的泪水；但既然泪水是悲伤的表示，他愿意你感受他人的痛苦如同自己的痛苦。正如你希望自己在那样痛苦中能得到帮助，你也应同样帮助他人，根据这句话：\BibleQuote{凡你们愿意人给你们做的，你们也要照样给人做\BibleRef{玛 7:12}}。因为与哭泣的人一同哭泣，却同时拒绝在可能时帮助那哭泣的人，是嘲弄而非虔诚的明证。总之，我们的救主与\Person{拉匝禄}的姐妹\Person{玛利亚}和\Person{玛尔大}一同哭泣，并以他泪水的见证显明了他内心无限的怜悯。但随后，作为真爱的证明，善工便随之而来：拉匝禄，那为之流泪的人，被复活并归还给他的姐妹。这才是真诚地与哭泣的人一同哭泣：消除了哭泣的原因。但你会说，他这样做是因为他有能力。是的，并未要求你做不可能的事：凡尽其所能的人，便已尽了全部责任。\par

% structure: p0034 source=h3 target=subsection reason=relative-heading-rank; base=h2

\subsection{第七章}

但是（正如我们开始评论的那样），基督徒仅仅远离邪恶是不够的，除非他也履行善工所包含的责任，这一点被主的那句话非常清楚地证明了，主警告说，那些虽然没有作恶、却没有行一切善的人，将被判处永火，祂说：\Quote{那时君王要对那些在他右边的人说：你们这些被诅咒的，离开我，到那为魔鬼和他的天使所预备的永火里去吧！因为我饿了，你们没有给我吃；我渴了，你们没有给我喝。}（见：\BibleRef{玛 25:41}）以及随后的话。祂并没有说：\Quote{你们这些被诅咒的，离开我，因为你们犯了杀人、奸淫或偷盗的罪。}因为他们被定罪、被判受永罚的地狱的惩罚，并不是因为他们作了恶，而是因为他们没有行善；并不是因为他们做了禁止的事，而是因为他们不愿意做那些命令的事。由此可以看出，那些甚至做了某些禁止之事的人，还能有什么希望呢？因为连那些仅仅没有做命令之事的人，都被判入永火。因为我不希望你这样自欺——如果你没有做某些事，就因为你做了某些别的事，因为经上记载：\Quote{凡遵守全部法律，却在一条上失足的，就成了违犯全部法律的。}（见：\BibleRef{雅 2:10}）因为亚当犯了一次罪，就死了；而你常常做那种别人只做一次就致死的事，你以为你能活吗？或者你认为他犯了大罪，因此被公正地判了更重的刑罚？那么，让我们想想他究竟做了什么。他违背了诫命，吃了树上的果子。那么，天主难道因为一棵树的果子就用死亡来惩罚人吗？不：不是因为那棵树的果子，而是因为对诫命的轻视。因此，问题不在于过犯的性质，而在于对诫命的违犯。那曾命令亚当不要吃树上果子的那一位，也命令你：不要毁谤，不要说谎，不要诽谤，不要听信诽谤者，完全不发誓，不要贪图，不要嫉妒，不要醉酒，不要贪饕，不要以恶报恶于任何人，要爱你们的仇敌，祝福那诅咒你们的，为那辱骂并迫害你们的祈祷，有人打你的右脸，把另一面也转给他，不在世俗法庭诉讼，这样，如果有人想夺走你的财物，你要欣然失去它们，要逃离贪婪，提防一切骄傲和自夸的罪，并效法基督的榜样，生活谦卑温良，完全避免与恶人交往，甚至不与行淫的人、贪婪的人、毁谤他人的人、嫉妒的人、诽谤者、醉酒的人、勒索的人一同吃饭。现在，如果你在任何这样的事上藐视祂，那么，如果祂饶恕了亚当，祂也会饶恕你。是的，祂可能更有理由饶恕你，因为亚当仍然无知无经验，并且没有被任何先前犯罪并因罪而死的人的榜样所约束。而你拥有如此多的榜样，在律法之后、先知之后、福音之后、宗徒之后，如果你仍然决心犯法，我看不出还能怎样赦免你。\par

% structure: p0036 source=h3 target=subsection reason=relative-heading-rank; base=h2

\subsection{第八章}

你因童贞的德行而自夸吗？要知道亚当和厄娃堕落时也是童贞的，当他们犯罪时，他们身体的完美纯洁并没有帮助他们。犯罪的童贞女应被比作厄娃，而非玛利亚。我们不否认，在今世有补赎的补救，但我们提醒你更应盼望赏报，而非寻求宽恕。因为那些期望童贞华冠的人竟然求取宽恕，而那些甚至已断绝合法之事的人竟然做出不合法的事，这是可耻的；因为必须记住，缔结婚姻的盟约是合法的。正如那些因爱基督和天国荣耀而轻视婚约的人应受赞扬，同样那些因不节制的快乐而在将自己奉献给天主后又求助于宗徒的补救的人应受谴责。因此，正如我们所说，拒绝婚姻的人轻视的不是不合法的事，而是合法的事。如果那类人发誓，如果他们诽谤他人，如果他们诋毁人，或者他们耐心听信诋毁者，如果以恶报恶，如果他们贪图他人的财产，或对自己的财产贪得无厌，如果他们怀有报复或嫉妒的毒害，如果他们说出或思想任何不合乎法律或宗徒的规条，如果为了取悦肉体而打扮装饰自己，如果他们做出任何其他不合法的事（这是常见的事），那么他们轻忽了合法的事，却实行不合法的事，这对他们有什么益处呢？如果你希望轻视合法的事对你有益，就要小心不做任何不合法的事。因为害怕本性上较小的，而不害怕本性上较大的，是愚蠢的[或因不避免那些被禁止的事，而当允许的事遭到轻视]。因为宗徒说：\BibleQuote{未婚的女子挂虑主的事，怎样使天主喜悦，为能在身体和灵魂上都成为圣的；但已婚的妇女挂虑世俗的事，怎样使丈夫喜悦。}\BibleRef{格前 7:34}他以此证明，已婚妇女因思想世俗的事而取悦丈夫，未婚女子则因没有挂虑世俗的事而取悦天主。那么，请告诉我，那个没有丈夫却挂虑世俗的事的人，究竟想取悦谁？在这种情况下，已婚妇女不是胜过她吗？是的，因为已婚妇女因挂虑世俗的事至少取悦了丈夫，而另一个既没有丈夫可取悦，也不能取悦天主。但我们不应缄默他所说的：\BibleQuote{未婚的女子挂虑主的事，怎样使天主喜悦，为能在身体和灵魂上都成为圣的}[他所说的，\Quote{她关心的是主的事}；她不关心世界或人的事，而是天主的事]。那么，什么是主的事呢？让宗徒告诉我们：\BibleQuote{凡是圣洁的，凡是公义的，凡是可爱的，凡是有良好声誉的，若有什么德行，若有什么学说的赞美}：这些就是主的事，是圣洁且真正宗徒的童贞女日夜思想、毫不停息地默想的。主的事是死人的复活，主的事是不朽，主的事是不腐，主的事是福音中所应许给圣人的太阳的光辉，正如福音所载：\BibleQuote{那时，义人将在他们父的国里发光如同太阳}\BibleRef{玛 13:43}。主的事是义人在天上许多住处，主的事是所结的果实，有三十倍、六十倍、一百倍的。那些思想这些事并思想通过什么行为才能配得这些事的童贞女，所想的就是主的事。主的事也是新约和旧约的法律，在其中闪耀着他口中的圣言；如果有童贞女不停默想这些事，她们所想的正是主的事。那时，先知的话就在她们身上应验了：\BibleQuote{永恒的根基是在坚固的磐石上，天主的诫命在圣善妇人的心中。}\BibleRef{训 26:24}\par

% structure: p0038 source=h3 target=subsection reason=relative-heading-rank; base=h2

\subsection{第九章}

接下来的语句是\Quote{怎样使天主喜悦}——我说，是天主，不是人——\Quote{为能在身体和灵魂上都成为圣的。}他不是说仅在某一肢体或身体上成圣，而是说在身体和灵魂上成圣。因为肢体只是身体的一部分，而身体是所有肢体的结合。因此，当他说她应在身体上成圣时，他证明她应当在所有肢体上成圣，因为如果在一个肢体中仍残留着败坏，其他肢体的圣化将是无效的。同样，\Emph{她}也不会在身体上（身体由所有肢体组成）成圣，如果她被其中一个肢体的污染所玷污。但为了使我所说的更加明显和清楚，假设一个妇女通过所有其他肢体的圣化而洁净，但只用她的舌头犯罪，即诽谤他人或作伪证，那么她所有其他肢体能否保证这一个肢体无罪，还是所有其他肢体都因这一个肢体而受审判？因此，如果其他肢体的圣化无效，即使只有一个肢体有过失，那么若所有肢体都被各种罪的罪愆所败坏，一个肢体的完美又有什么用呢？\par

% structure: p0040 source=h3 target=subsection reason=relative-heading-rank; base=h2

\subsection{第十章}

\Emph{因此}，我恳求你，贞女啊，不要单因你的纯洁而自夸，也不要信赖某一肢体的完美；而要依照宗徒的教导，保持你全身的圣洁。洁净你的头，除去一切污秽，因为头在领受祝圣之油后，若再被番红花汁或粉末，或任何其他颜料玷污，或用金子、宝石或任何其他世俗装饰来装扮，便是可耻的；因为它已因天上的装饰而发光。无疑地，将任何世俗和尘世的装饰置于天主的恩宠之上，是对天主恩宠的严重侮辱。其次，洁净你的额头，使它为人而非为天主的事工而脸红，并显出那种羞愧——它不是导致罪恶，而是带来天主的恩宠，正如圣经所载：\Quote{有一种羞愧导致罪恶，也有一种羞愧带来天主的恩宠（见：\BibleRef{训 4:21}）。}也要洁净你的颈项，使它不致将你的发丝用金网和项链缠绕，而应佩戴圣经所说的装饰：\Quote{不要让\BibleRef{箴 3:3}仁慈和信德离开你，}并将它们挂在你的心上，如同挂在你的颈上。洁净你的眼睛，使它们远离一切私欲偏情，永不转离对穷人的注视，并保持它们在天主所造的那种纯洁中，远离一切脂粉。洁净你的舌头，远离谎言，因为\Quote{说谎的\BibleRef{智 1:11}口舌毁灭灵魂}；洁净它远离毁谤、发誓和伪誓。我请求你不要以为我先说舌头应从发誓中洗净，后说从伪誓中洗净是次序颠倒；因为人若根本不起誓，就更易避免伪誓，如此便可在他身上应验那句话：\Quote{保守你的舌头远离邪恶，你的嘴唇不说诡诈的话。}要记住宗徒所说：\Quote{祝福，\BibleRef{罗 12:14}不要诅咒。}但请常常记起以下的话：\Quote{你们要小心，不可有人以恶报恶，以咒骂还咒骂，反而要祝福，因为你们蒙召正是为此，好使你们承受祝福。}以及另一处：\Quote{若有人在\BibleRef{雅 3:2}言语上不犯罪，他便是一个完人。}因为你用以宣认天主、向他祈祷、祝福他、赞美他的嘴唇，若被任何罪恶的污秽所玷污，那是可耻的。我不知道一个人怎能用那说谎、毁谤或中伤他人的舌头来向天主祈祷。天主垂听圣洁的嘴唇，并迅速答复那些由无玷之舌所发出的祈祷。也要洁净你的耳朵，使它们只听从圣洁和真实的话语，永不接纳淫秽、可耻或世俗的言辞，也不容忍任何人中伤他人，因为经上记载：\Quote{用荆棘\BibleRef{训 28:24}堵塞你的耳朵，不要听邪恶的舌头，}好使你有分于那被称为\Quote{在\BibleRef{伯后 2:8}听觉和视觉上都是义人}的人，即他既不以眼也不以耳犯罪。也要洁净你的双手，\Quote{使它们不伸出来接受，却紧闭着不给予，}并\BibleRef{训 4:31}使它们不急于打击，却随时准备施行一切仁慈和虔敬的事工。最后，洁净你的双脚，使它们不跟随那引向盛大而奢侈的世俗宴会的宽阔大道，而要走那引向天堂的崎岖窄路，因为经上记载：\Quote{为你的脚\BibleRef{箴 4:26}修直道路。}要承认，你的肢体是由天主造主为你所造，不是为了恶习，而是为了德行；当你洁净了全身的肢体，除去一切罪恶的玷污，它们在整个身体中都成为圣洁时，那时你才明白这贞洁于你有益，并满怀信心地期盼贞洁的赏报。\par

% structure: p0042 source=h3 target=subsection reason=relative-heading-rank; base=h2

\subsection{第十一章}

\Emph{我相信}，我现在已经简要但同时也充分地阐述了妇女的\OriginalTerm{身体的纯洁}{purity of body}所包含的意义：接下来我们应当学习什么是\OriginalTerm{精神纯洁}{pure in spirit}，即一个人不能做的事，甚至也不能在思想中想象。因为她在身体和灵魂上都是圣洁的，凡不在心思或心中犯罪的人，知道天主也鉴察人心；因此，她竭尽全力保持心思和身体都无罪。这样的人知道经上记载说：\BibleQuote{你要谨守你的心，\BibleRef{箴 4:23} 胜过一切}；又说：\Quote{天主喜爱圣洁的心，一切无玷污的人都为祂所悦纳}；另有一处说：\BibleQuote{心里洁净的人是有福的，\BibleRef{玛 5:8} 因为他们要看见天主。}我认为这最后的话是针对那些良心不指控他们犯任何罪的人说的；关于他们，我认为若望在他的书信中也说过：\BibleQuote{如果我们的心\BibleRef{若一 3:21}不责备我们，那么我们在天主前就放心，无论求什么，必从祂得着。}我不愿你以为自己逃脱了罪的控诉，尽管欲望没有付诸行动，因为经上记载：\BibleQuote{凡注视妇女，\BibleRef{玛 5:28} 有意贪恋她的，他已在心里与她犯了奸淫。}也不要这样说：\Quote{我确实起了念头，但没有付诸实行}；因为凡不能做的事，连想都是非法的。因此，圣伯多禄也颁布了这样的训令：\BibleQuote{你们要净化\BibleRef{伯前 1:22}自己的灵魂}；如果他不知道有\OriginalTerm{灵魂的玷污}{defilement of the soul}这回事，他就不会表达净化灵魂的愿望。但我们还应非常仔细地思考那段话：\BibleQuote{这些人\BibleRef{默 14:4}是那些没有与妇女玷污自己的人，因为他们保持了\OriginalTerm{童贞者}{virgins}，他们跟随羔羊，无论祂到哪里去}；并要反思，如果这些人加入了天主的随从队伍，穿越诸天各处，仅凭\OriginalTerm{贞洁}{chastity}和纯洁的功德，那么是否还有其他途径，使童贞得以借助这些途径达到如此伟大真福的荣耀。但我们从哪里知道这一点呢？从以下段落（如果我没记错的话）中记载着：\Quote{这些人是从人间买来的，作为初熟的果子归于天主和羔羊，在他们口中没有谎言，因为他们在天主宝座前是没有瑕疵的。}那么你看，他们被称为紧紧跟随主的脚步，不是仅凭一个肢体，而是那些除了童贞之外，还过着远离一切罪恶污染生活的人。因此，尤其要让童贞女轻视婚姻，因为这样她比其他人更安全，更容易完成已婚者也需要做到的事：即远离一切罪恶，遵守法律的一切诫命。因为如果她不结婚，却沉溺于即便是已婚妇女也被命令远离的事情，那么不结婚又有什么益处呢？虽然任何基督徒都不被允许犯罪，凡通过\OriginalTerm{圣洗圣事}{spiritual bath}的\OriginalTerm{圣化}{sanctification}而得洁净的人，都理应过无玷的生活，以便与那描述为\BibleQuote{没有\BibleRef{弗 5:27}斑点、皱纹或任何这类毛病的}教会完全相符；但童贞女更应达到这一标准，因为既没有丈夫、儿子，也没有任何其他必需事物阻碍她完全实行圣经的要求；如果她失败，也不能用任何借口为自己辩护。\par

% structure: p0044 source=h3 target=subsection reason=relative-heading-rank; base=h2

\subsection{第十二章}

童贞女啊，持守你的志向，这志向注定要得丰厚赏报。在主面前，贞洁与纯洁之德是卓越的，只要它不被其他种类的犯罪与邪恶所玷污。认清你的身份，认清你的地位，认清你的志向。你被称为基督的新娘；注意不要做出任何与你所宣称要许配的那一位不相称的行为。他若在你身上察觉一丝不忠，便会迅速写下休书。同样，凡接受那些作为人类订婚定金而赐予的礼物者，立刻便开始殷勤而勤奋地向仆人、亲近之人和朋友询问：那青年人的品性如何，他特别喜爱什么，他接受什么，他过着怎样的生活，他有什么习惯，他耽于何种享乐，他在哪些追求中找到最大的快乐和喜悦。当她得知这些事之后，便事事如此行事，使她的服事、她的喜乐、她的勤勉以及她整个生活方式，都能与未婚夫的品性相合。而你，以基督为你的新郎，要从你的新郎的仆人和亲近之人那里询问他的品性；是的，要热心而巧妙地问：他特别喜爱什么，他喜欢你的衣着有何种安排，他想要何种装饰。让他的最亲密的同伴\Person{伯多禄}告诉你，他甚至在书信中也不允许已婚妇女妆扮自身，如他所写：\BibleQuote{同样，你们做妻子的，要服从自己的丈夫，好使那些不信从主话的，为了妻子无言的品行而受感化，因为他们看见了你们怀有敬畏的贞洁品行。你们的装饰不应是外面的发型、金饰、或衣服的华丽，而应是那藏于内心，从不朽的温柔宁静的精神而来的装饰，这在天主前是极宝贵的。}让另一位宗徒也告诉你，真福\Person{保禄}，他写信给\Person{弟茂德}时，赞同关于信德妇女行为的同样道理：\BibleQuote{同样，也愿女人以端庄和持守的服饰来装饰自己，不靠卷发、金饰、珍珠或华贵的衣服，而要以善行和端正的行为，这才适合自称为虔敬天主的女人。}\par

% structure: p0046 source=h3 target=subsection reason=relative-heading-rank; base=h2

\subsection{第十三章}

但也许你会说：\Quote{为什么宗徒们没有将这些事吩咐给贞女们？}因为他们认为没有必要，唯恐这样的劝勉给予她们，反而更像是侮辱而不是建树。事实上，他们也不会相信贞女们竟会大胆到如此地步，将连已婚妇女都不被允许的属肉体和属世的装饰据为己有。毫无疑问，贞女应当装饰和打扮自己；因为如果她不以整洁和美丽的形式出现，怎能取悦她的未婚夫呢？她务必被装饰，但她的装饰应是内在和属灵的，而非属肉体的；因为天主渴望在她身上的不是身体的美丽，而是灵魂的美丽。因此，你既然渴望你的灵魂被天主所爱并居住其中，就当殷勤地装饰它，以属灵的衣裳打扮它。不要在它身上看到任何不雅或可厌之物。让它以公义的金子发光，以圣洁的宝石闪烁，以极宝贵的纯洁珍珠灿烂；不要穿细麻和丝绸，而要穿上慈悲和虔诚的外衣，正如经上所记：\BibleQuote{穿上怜悯、慈爱、谦逊，等等}\BibleRef{哥 3:12}也不要让贞女追求脂粉或任何颜料的美丽，而要有纯真和质朴的光辉，端庄的玫瑰红，以及可敬的羞赧的紫色光辉。让她以天上教理的碱水洗涤，并以一切属灵的沐浴洁净。不要在她身上留下任何恶意或罪污。并且，唯恐她有时散发出罪恶的恶臭，让她满被智慧和知识的最甜美的香膏所浸透。\par

% structure: p0048 source=h3 target=subsection reason=relative-heading-rank; base=h2

\subsection{第十四章}

天主寻求这种装饰，渴望一个如此装扮的灵魂。请记住，你被称为天主的女儿，正如祂所说：\Quote{女儿，请听，请留意。}而你自己，每当称天主为父时，就证明你是天主的女儿。因此，如果你是天主的女儿，要小心不要做任何不配于你父天主的事；反而要身为天主的女儿，做一切事。思考一下世上贵族的女儿们如何行事，她们习惯于何种习性，通过何种训练来陶冶自己。她们中的一些人，如此谦逊、如此庄重、如此自制，以致在人的高贵方面超越了其他人的习性；为了不让自己的失败给可敬的父母带来任何耻辱，她们努力通过自己在世上的行为方式为自己获得另一种本性。因此，你也要留意你的根源，考虑你的出身，注意你高贵的光荣。承认你不仅是人的女儿，而且是天主的女儿，并饰以神圣出生的高贵。如此向世界呈现你自己，使你天上的出生在你身上被看见，你神圣的高贵清晰闪耀。愿在你身上有新的尊严、令人赞叹的美德、显著的谦逊、奇妙的忍耐、合于贞女的步伐，带着真正羞怯的态度，言语总是谦逊，只在适当的时候说出，以致凡看见你的人都会惊叹地说：\Quote{在人中间，这种新尊严的表现是什么？这种引人注目的谦逊是什么？这种平衡的美德是什么？这种成熟的智慧是什么？这不是人类训练的成果，也不是单纯的人类纪律的成果。某种天上的事物，在那尘世的躯体中，向我散发芬芳。我真的相信天主确实居住在有些人身上。}当他得知你是基督的侍女时，他将更加惊叹，并思考主人是多么奇妙，因为他的侍女表现出如此的美德。\par

% structure: p0050 source=h3 target=subsection reason=relative-heading-rank; base=h2

\subsection{第十五章}

因此，你如果愿意与基督同在，就必须按照基督的表样生活，祂远离一切邪恶和恶毒，甚至不报复仇敌，反而为他们祈祷。因为我不希望你视那些灵魂为基督徒，他们（我不说）恨自己的兄弟或姐妹，而是不在天主作证下，全心、全良心爱近人；因为基督徒有责任，效法基督本身，甚至爱仇敌。如果你想与圣人们共融，就要从恶意和罪恶的思想中洁净你的心。不要让任何人欺骗你，不要用诱骗的言语迷惑你。天国的宫廷只接纳圣洁、正义、纯朴、无辜和纯洁的人。邪恶在天主面前没有位置。凡愿意与基督一起为王的人，必须远离一切邪恶和诡诈。没有比仇恨任何人、想伤害任何人更令天主厌恶和可憎的事；而爱所有的人最令祂喜悦。先知知道这一点，当他教导时说：\Quote{你们爱慕主的，要憎恨邪恶。}\par

% structure: p0052 source=h3 target=subsection reason=relative-heading-rank; base=h2

\subsection{第十六章}

注意，不要在任何方面喜爱人的光荣，免得你的份也被算在那些享受光荣的人中，对他们曾说：\Quote{你们既然彼此求光荣，\BibleRef{若 5:44}你们怎能相信？}并且论到他们，先知曾说：\Quote{增加他们的邪恶，增加地上自夸者的邪恶}；在别处又说：\Quote{你们因自己的夸耀而蒙羞，因在主面前的责斥而蒙羞。}因为我不希望你顾念那些世上的贞女，而非基督的贞女；她们忘记了自己的意向和誓愿，喜乐于美味，为财富而喜悦，并以纯属肉体的高贵出身而夸耀；如果她们确实相信自己天主的女儿，就决不会在神圣的起源之后，仰慕单纯的人类高贵，也不会以任何荣誉的尘世父亲为荣；如果她们感到天主是她们的父，就不会喜爱任何与肉体有关的高贵。愚昧的女人，你为何因你出身的尊贵而自夸，并以此为乐？天主起初创造了两个人，整个人类种族都由他们繁衍而来；因此，尘世的高贵并非出于天性的公平，而是出于邪恶欲望的野心。无疑，我们都因神圣洗礼的恩宠而成为平等的；在第二次诞生所生的人中，不能有差异；藉此，富人和穷人、自由人和奴隶、贵族和低贱的人，都成为天主的儿子。因此，单纯尘世的等级被天上荣耀的光辉所遮蔽，从此不再被记；那些以前在世俗荣誉上不平等的人，如今在天上神圣的高贵荣耀中平等地被装饰。在他们中间，没有低贱出生的余地；也没有人比另一个被神圣出生的尊威所装饰的人更低微；除非在那些不认为天上的事物应当优先于地上事物的人看来。他们中间不能有世俗的夸耀，如果他们反思：在较小的事情上，他们将自己置于那些他们知道在较大事情上与自己平等的人之上，这是多么虚空；并且，他们将以在地上低于自己的人，视为那些他们相信在关乎天上的事上与自己平等的人。但你，既是基督的贞女，而非世界的贞女，要逃避现世生命的一切光荣，以便获得来世所应许的光荣。\par

% structure: p0054 source=h3 target=subsection reason=relative-heading-rank; base=h2

\subsection{第十七章}

避免争辩的言辞和敌意的起因；也要逃避一切不和睦和争斗的场合。因为，如果按宗徒的教训，\BibleQuote{主的仆人\BibleRef{弟后 2:24}不应争辩}，那么，主的婢女更当如此，她的心意应更温柔，正如她的性别更羞怯和退让。约束你的舌头不说坏话，把法律的缰绳套在你的口上；这样，你若要说话，就只当沉默是罪时才说话。当心不要说出任何可能被合理指责的话。一言既出，如石已投：因此说出口前应深思良久。有福的确是有福的嘴唇，从不说出他们希望收回的话。贞洁心灵的谈论自身也应当是贞洁的，总是更能建立而不是伤害听者，按照宗徒的命令，他说：\BibleQuote{任何\BibleRef{弗 4:29}腐败的言语不可从你们的口中发出，只可说那为建立信德有益的话，好能将恩宠带给听者。}在天主眼中宝贵的是那舌头，除了谈论神圣之事不知如何措辞；圣洁的是那嘴巴，不断涌出天上的话语。用圣经的权威贬斥那些不在场之人的诽谤者，视他们为心术不正之人，因为先知也将这一点列为完人的美德之一：如果在义人面前一个心术不正的人提出无法证实的邻人恶事，他便被贬为无有。因为你不可耐心地听别人受诽谤，你也不愿别人对你这样做。的确，一切违反基督福音的都是不义的，如果你容忍别人遭受任何你自身会感到痛苦的对待，也是如此。习惯让你的舌头总是谈论善人，让你的耳朵更乐于听善人的赞美而非恶人的定罪。注意你的一切善行都是为天主的缘故而行，知道每一件善行你将只得到报酬，只要你出于敬畏和爱他而做。努力做圣洁而非显得圣洁，因为被认为是你所不是的毫无益处；当你没有别人以为你有的，并假装拥有你没有的，你就犯了双重的罪。\par

% structure: p0056 source=h3 target=subsection reason=relative-heading-rank; base=h2

\subsection{第十八章}

你更应喜爱斋戒胜过宴乐，牢记那位没有离开圣殿、日夜以斋戒和祈祷事奉天主的寡妇。现在，如果她是一位寡妇，一位犹太寡妇，尚且如此，那么基督的贞女现在应该达到什么地步呢？你当爱慕神圣言语的盛宴胜过一切，渴望充满属灵的美味，寻求那滋养灵魂的食物，胜过只愉悦身体的。逃避一切肉和酒，作为引发情欲的热源和挑逗。只有当胃部不适或身体极度虚弱需要时，才使用一点酒。制服愤怒，抑制敌意，凡是做后会引起懊悔的事，都要当作引致立即之罪的憎恶而避免。那渴望成为天主居所的心灵应当十分平静安宁，不受一切愤怒骚动的干扰，正如祂借先知所证实的说：\BibleQuote{\BibleRef{依 66:2}我将安息在谁身上，除非那谦卑安静、对我的话战栗的人？}相信天主是你一切行为和思想的见证，当心不要做或想任何有损于神圣视线的事。当你渴望祈祷时，要表现出与主交谈者应有的心态。\par

% structure: p0058 source=h3 target=subsection reason=relative-heading-rank; base=h2

\subsection{第十九章}

当你重复一篇圣咏时，要思考你重复的是谁的话，并以灵魂真正的痛悔为乐，胜过以颤音的悦耳为乐。因为天主更看重赞美祂之人的眼泪，而非他声音的美妙；如先知所言：\Quote{应以敬畏事奉主，且以颤抖欢喜。}现在，哪里有恐惧和颤抖，哪里就没有高声，而是谦卑的心，伴随哀叹和泪水。在一切事上展现勤奋；因为经上记载：\BibleQuote{以疏忽从事主工作的人，是可咒诅的。}\BibleRef{耶 48:10}让恩宠随年岁增长，让正义与年龄俱增；让你的信德随着你越年长而显得越完善；因为耶稣给我们留下了如何生活的榜样，祂不仅在身体方面年岁增长，也在智慧和属神的恩宠上，在天主和人面前增长。要将你没有察觉自己变得更好的所有时间视为确已失去。坚守你所立下的贞洁心愿直至最后；因为美德不仅在于开始，更在于完成，正如主在福音中所说：\BibleQuote{坚忍到底的，必要得救。}\BibleRef{玛 10:22}因此，要谨慎，免得你给任何人甚至邪恶欲望的借口，因为你的天主已与你订婚，祂是忌邪的；因为对基督犯奸淫比对自己的丈夫犯奸淫更有罪。所以，你要成为所有人的生活模范；成为榜样；并在实际行为上超越那些先你奉献于贞洁的人。在各方面显示你是一个贞女；不让任何腐败的污点被指控到你身上。让身体在纯洁上完美的人在行为上也无可指责。正如我们在信的开头所说，你已成为献于天主的祭品，这样的祭品无疑也将自己的圣德赋予他人，如此，每个人都相称地领受，他自己也能成为圣化的分享者；那么，就让其他贞女也通过你，如同通过一种神圣的奉献，得到圣化。你要在一切事上向他们显示圣洁，以致凡接触你生活的人，无论是听见还是看见，都能感受到圣化的力量，并感到如此多的恩宠从你的行为方式传递给他，以致他渴望效仿你，自己也就成为奉献于天主的祭品。\par

% structure: p0060 source=h2 target=section reason=relative-heading-rank; base=h2

\section{第三封信}

塞维鲁致圣保禄主教书\par

得知你所有的厨师都已离开你的厨房（我相信是因为他们对不得不为廉价豆类菜肴履行职责感到愤慨）后，我从我们自己的作坊里派了一个小男孩给你。他非常善于煮淡色的豆子，用醋和酱汁腌制朴素的甜菜根，也能为饥饿的隐修士们准备廉价的粥。但关于胡椒或拉瑟草，他一无所知；而他精通小茴香，尤其擅长用香甜的植物捣动研钵（发出吵闹声）。他有一个缺点：任何花园都不允许他进入；因为如果被放进去，他会用剑砍倒他所能触及的一切，而且绝不会满足于仅仅屠杀锦葵。然而，在搜集燃料方面，他不会欺骗你：他会烧掉任何挡道的东西；他会砍倒并毫不犹豫地动手破坏房屋，拿走家里的旧梁木。那么，我们把他连同这些性格和美德赠予你；我们希望你把他看作一个儿子，而非仆人，因为你不羞于做极小生灵的父亲。我本人本希望代他服侍你；但如果善意在某种程度上可以代替行动，请只回报我，在你的早餐和愉快的晚餐中记得我，因为做你的奴隶比做他人的主人更合适。请为我祈祷。\par

% structure: p0063 source=h2 target=section reason=relative-heading-rank; base=h2

\section{第四封信}

致同一人，论其智慧与温和\par

我们神圣宗教的忠实解说者如此安排一切，以致将来不再为犯过者留有余地：例如，你以如此伟大的圣德品格向我们应许了什么，无非是放下一切错误，过蒙福的生活？在这方面，我看到最大的赞美配得上你的美德，因为你甚至通过劝诫改变了未受教化的心灵，并将其吸引到卓越的状态。但如果你只是通过向受过教育的心灵灌输智慧而强化它们，就不会显得如此奇妙；因为聪明人对虔诚有一种亲近，而粗野的本性不容易被争取到严格的一边。正如那些用石头雕刻动物形状的人，当他们用凿子击打非常坚硬的岩石时，从事的是相当困难的工作，而那些尝试较软材料的人感到手因易于加工而受助，并且工人的劳作在困难时应当被给予最高荣誉；同样，阁下，应单独给予你称赞，因为你使未经琢磨和粗野的心灵，从罪的黑暗中被释放，既思考人性的事，又理解神圣的事。\par

色诺克拉底，远为最博学的哲学家，同样受人尊崇。他藉严厉的劝诫，成功使奢侈被战胜。因为有一次，某个名叫\Person{波勒莫}的人，酩酊大醉，在听众纷纷涌向\Person{色诺克拉底}学园时，公开从一场夜间狂欢中踉跄而出，他也进入该处，无耻地坐在一群门徒中间，穿着他从宴席中出来的那身衣服。花冠覆盖着他的头，他却并不因自己看上去与众不同而感到羞耻，因为实际上，长时间痛饮已扰乱了他的理智，而理智本是理性的所在。在场的其他人开始强烈抱怨，因为如此不合适的听众竟混入一群文人之中，但大师本人却毫不在意，反而开始论述道德科学和节制的律法。教师的影响力如此强大，以至于那个无耻的闯入者的心被说服去爱慕谦逊。于是，\Person{波勒莫}首先十分窘迫地摘下头上的花冠，承认自己是门徒。后来，他逐渐使自己完全符合尊严所要求的义务，并彻底献身于谦逊的表现，以至于品格的荣耀改变掩盖了他以前生活的习惯。我们钦佩您的教导中的这一点：无需使用任何威胁，也无需诉诸任何恐怖，您就将迷途的心灵引向对天主的敬拜；因此，即使是一个秩序不佳的理智，也应相信与其与少数人持有不义的见解，不如与众人一起过美好幸福的生活。\par

% structure: p0067 source=h2 target=section reason=relative-heading-rank; base=h2

\section{第五封书信}

致某位不知名人士，恳求他温和对待其兄弟。\par

虽然我的大人和兄弟已经恳求您的尊贵，确保\Person{图图斯}最安全，但允许我在信中推荐同一个人，以便通过加倍请求，使他更加安全。因为容我们假定，青春的过错和尚未定型的年龄的错误伤害了他，给他的早年留下了污点；然而，一个尚不知道正当行为是什么的人，几乎无可指责地犯了错。因为当他达到正确的思想和反省时，他通过更好的思考明白了戏剧生活应受谴责。然而，他不能完全摆脱他的过错，除非他藉天主的助佑洗去其罪疚，因为通过天主教所获得的补救，他改变了观点，否认自己享受一个不太光彩的场所，并使自己从民众的目光中退隐。\par

\Emph{论上述大师}\par

因此，既然神律和国法都不允许忠信的身体和圣化的心灵展示可耻但令人愉悦的表演，并展示庸俗的娱乐手段，尤其因为如果在某种程度上对自身的贞洁奉献造成伤害，万一任何因圣洗圣事而更新的人退回到他旧日的放荡，那么，您的尊贵理应眷顾良好的意向，以使那因天主的良善而开始虔诚义务的人，不致被迫陷入戏院的陷阱。然而，如果您们命令他采取其他适合我们共同国家需要的行动，他并不拒绝对您们所有人的判断表示顺服。\par

% structure: p0072 source=h2 target=section reason=relative-heading-rank; base=h2

\section{第六封书信}

致\Person{萨尔维乌斯}：投诉乡民受骚扰、财产遭掠夺\par

辩论的激情理应在法庭业务期间达到白热化；因为日复一日奋斗的勤劳之臂应展现有力的动作。但当宏亮的雄辩已吹响退兵号，并退到和平的树林和宜人的住所时，理应将闲散的怨言搁置一旁，停止发出无效的威胁。因为我们知道，得胜的骏马在退出竞技场后，便在厩中极其安静地休息。持续的恐惧和不确定的胜利棕榈不再困扰它们，而是终于系在和平的食槽边，不再敬畏驱策它们的主人，享受着先前盛行的那种不安竞争的甜美遗忘。同样，让那夸耀的士兵在服役期满后，乐于悬挂战利品，并耐心承受年岁的重担。\par

但我并不完全理解您为何乐于恐吓可怜的农夫；我也不明白您为何想用缺乏生计的恐惧来骚扰我的乡下人；仿佛我真不知道如何安慰他们，消除他们的恐惧，并向他们表明没有像您所说的那样可怕的缘由。我承认，当我们忙于平地时，我常被您雄辩的武器吓到，但我也常常尽力回击。我当然与您一同学习了，农夫们根据什么权利和顺序被要求归还，谁有资格进行法律程序，谁没有资格进行法律程序。您说\Person{沃卢西亚人}希望您被带回来，并且您常常在愤怒中重复，您将把我的乡下人从我的小领地中撤走；而您本人，正如我所希望和渴望的，受古老亲属关系的纽带约束于我，现在却鲁莽地威胁，要抛弃我们的协议，抓住我的人。我请问您卓越的知识，是否有一种法律为律师，另一种为私人，是否在罗马是一回事，而在\Place{马塔鲁姆}完全另一回事。\par

在此期间，我不知道你曾是\Place{沃卢西安}产业的主人，因为\Person{狄约尼削}被认为是保留了该产业的所有权，而他从不缺少继承人；他生前惯于用粗鄙言语的毒舌嘲弄许多人。当时有一个\Person{波菲利}，\Person{齐柏林努斯}之子，但他并非名副其实地被称为\Person{齐柏林努斯}之子。他通过军役隐瞒了出身的疑问，并且为了驱散额上的阴云，他参与了殷勤的服务和自愿的服从。他常与我在一起，无论是在家还是在广场，经常请我作为他的辩护人向我父亲陈情，并在法官面前作他的辩护律师。有时我甚至劝阻\Person{狄约尼削}，觉得他不该为了二十亩田而对\Person{波菲利}发泄粗鄙的辱骂。\par

看，这就是你那卓越的审慎威胁我的代理人的原因，尽管你不是该地的主人，却到处提到我的庄户人。但如果你自称是\Person{波菲利}的继承人，你必须知道二十亩的狭小空间确实无法由一位农夫经营；或者，如果你顾及自己应有的尊严并决心维持它，而不愿自称为\Person{波菲利}的继承人，那么很明显，有权这样做的人可以提起诉讼，以便与那些在该地上没有财产的人对簿公堂。但如果你仔细查考此事，你会发现追回该地的努力尤其落在我身上。因此，我极为尊敬的先生和兄弟，你理当保持和平，并重新与我为友，同时屈尊前来私下会谈。请你停止惊扰那些不活跃且易受惊吓的人，并在远处夸耀你的言辞。然而，请相信我，我对你的高昂精神感到欣喜，绝无冒犯；因为我们既不是性格严苛的人，也不缺乏学识。至少让\Person{马克西米努斯}使你温顺吧。\par

% structure: p0078 source=h2 target=section reason=relative-heading-rank; base=h2

\section{第七封信}

致一位未知者，请求赐函。\par

灵魂的信德和虔敬无疑仍然存在，但这应当通过书信的证据显明出来，以便透过这种相互的礼节能获得更多的情感增长。因为正如一块肥沃的田地，若疏于耕作，便不能结出丰硕的果实，那土壤的优良品质也会因休息而不工作的人的懒惰而丧失；同样，我认为，除非那些远方的人通过书信如同在场一般被探望，否则内心的爱与善意也会变得衰弱。\par

\SourceRecord{Translated by Alexander Roberts.；Nicene and Post-Nicene Fathers, Second Series；Vol. 11.；Edited by Philip Schaff and Henry Wace.；Buffalo, NY: Christian Literature Publishing Co.,；1894.；<http://www.newadvent.org/fathers/3504.htm>.}
